\subsection{Bimanual Exposure and Execution}
\label{sec:task}

We study bimanual exposure and execution under deformable occlusion.
The arms first establish visibility and access to an occluded target, then complete a prescribed operation and any required restoration.
Because exposure can change both the occluder configuration and target state, subsequent actions must respond to the resulting interaction conditions.
We learn from a finite, fixed demonstration set and evaluate complete-task success within a finite horizon using visual history.

\subsubsection{System, Demonstrations, and Policy}

A single policy coordinates both arms, whose roles may change during execution.
The joint robot-environment state $s_t\in\mathcal S$ includes the arms, target, deformable occluder, and their interaction conditions.
The joint command $u_t\in\mathcal A\subseteq\R^{d_u}$ controls both arms and their end effectors.
The observation $o_t=(I_t^{1:V},\mathbf{x}_t^{\mathrm{prop}})$ contains images from $V\ge1$ cameras and the proprioceptive state $\mathbf{x}_t^{\mathrm{prop}}$.
Over $T$ control steps, the interaction follows
\begin{align}
 s_0&\sim p_0,\nonumber\\
 s_{t+1}&\sim P(\cdot\mid s_t,u_t),\quad 0\le t<T,\nonumber\\
 o_t&\sim O(\cdot\mid s_t),\quad 0\le t\le T.
 \label{eq:interactionmodel}
\end{align}
where $p_0$ is the initial-state distribution and $P$ and $O$ are the transition and observation kernels.
The transition kernel captures coupled arm motion, occluder deformation, and target interaction.

The fixed demonstration set $\mathcal D=\{\zeta_i\}_{i=1}^{N}$ records observations and expert commands for the same task objective:
\begin{equation}
 \zeta_i=(o_0^i,u_0^{\mathrm E,i},o_1^i,\ldots,
 u_{T_i-1}^{\mathrm E,i},o_{T_i}^i).
 \label{eq:demonstrations}
\end{equation}
Here, $T_i$ is the length of trajectory $i$ and the superscript $\mathrm E$ denotes expert commands.
For the analysis, the deployed policy is deterministic and selects $u_t=\pi(\mathcal H_t)$ from the available history $\mathcal H_t$, which contains observations $o_{0:t}$, past commands $u_{0:t-1}$, and any runtime context available before the current decision.

\subsubsection{Learning and Task Objectives}
Let $\pi_{\mathrm E}$ denote the expert policy and let $\nu$ be the population distribution of uniformly sampled demonstration control steps and their histories.
This uses the same control-step weighting as the deployment average introduced below.
The population imitation risk is
\begin{equation}
 R_\nu(\pi)
 :=\E_{(t,h)\sim\nu}\norm{\pi(h)-\pi_{\mathrm E}(h)}^2.
 \label{eq:populationimitationrisk}
\end{equation}
The empirical objective in Section~\ref{sec:method} uses a masked $L_1$ surrogate for this expert-action discrepancy.

The measurable completion function $\Phi:\mathcal S^{T+1}\to\{0,1\}$ evaluates exposure, the prescribed target operation, and any required restoration independently of fixed-time stage boundaries or arm-specific success criteria.
Task-specific criteria are given in Section~\ref{sec:experimentalsetup}.
The completion indicator and success probability are
\begin{equation}
 \Psi=\Phi(s_{0:T}),\qquad
 J(\pi)=\Prob_\pi(\Psi=1).
 \label{eq:visualtask}
\end{equation}
Here, $\Prob_\pi$ is the trajectory law induced by the system and policy.
Learning reduces $R_\nu(\pi)$, whereas deployment seeks to increase $J(\pi)$.
The following analysis states conditions under which these objectives can be related.

\subsubsection{Visual Assessability}
Let $\Omega_{\mathrm{eval}}\subseteq\mathcal S^{T+1}\times\mathcal I^{T+1}$ contain the state and image histories considered under the fixed evaluation setup, where $\mathcal I$ is the space of multi-view images.

\begin{assumption}[Visual assessability]
\label{ass:visualassessability}
There exists a measurable visual decision function $\mathcal V:\mathcal I^{T+1}\to\{0,1\}$
such that, for every $(s_{0:T},I_{0:T}^{1:V})\in\Omega_{\mathrm{eval}}$,
\begin{equation}
 \Phi(s_{0:T})=\mathcal V(I_{0:T}^{1:V}).
 \label{eq:visassumption}
\end{equation}
\end{assumption}
Visual assessability concerns task completion rather than full-state observability.
The evaluation domain includes temporary occlusion but excludes outcomes that remain visually unobservable.
All policy executions and counterfactual comparisons below are restricted to $\Omega_{\mathrm{eval}}$ almost surely.
